% Generated draft: requires mathematical and source review.
\documentclass[11pt]{article}
\usepackage[margin=25mm]{geometry}
\usepackage{fontspec}
\setmainfont{Latin Modern Roman}
\usepackage{amsmath,amssymb,mathtools,unicode-math}
\setmathfont{Latin Modern Math}
\usepackage{enumitem,needspace,xurl,hyperref}
\setlength{\parindent}{0pt}
\setlength{\parskip}{5pt}
\setlength{\emergencystretch}{3em}
\newcommand{\N}{\mathbb{N}}
\newcommand{\Z}{\mathbb{Z}}
\newcommand{\Q}{\mathbb{Q}}
\newcommand{\R}{\mathbb{R}}
\newcommand{\C}{\mathbb{C}}
\newcommand{\F}{\mathbb{F}}
\newcommand{\A}{\mathbb{A}}
\newcommand{\PP}{\mathbb{P}}
\newcommand{\E}{\mathbb{E}}
\newcommand{\eps}{\varepsilon}
\newcommand{\dd}{\,\mathrm{d}}
\begin{document}
\section*{UnsolvedMath ID: 700373}
\textbf{Source problem code:} catalogue entry
\subsection{Definitions and mathematical statement}
Let \(G\) be a finite group with identity \(1\). An involution is an element of order exactly \(2\). A normal set of involutions is a subset \(D\subseteq G\) such that every element of \(D\) is an involution and \(gDg^{-1}=D\) for every \(g\in G\). Define the undirected simple graph \(\Gamma(D)\) by
\[
V(\Gamma(D))=D,\qquad E(\Gamma(D))=\bigl\{\{a,b\}:a,b\in D,\ a\ne b,\ ab=ba\bigr\}.
\]
This is the simple-graph interpretation of the supplied condition \(ab=ba\ne 1\): for involutions, \(ab\ne 1\) is equivalent to \(a\ne b\).

\textbf{Classification problem (A.\ A.\ Makhn\"ev, 1978).} Describe the finite groups \(G\) with a normal set of involutions \(D\) for which \(\Gamma(D)\) is disconnected. For nonempty \(D\), disconnectedness means that there exist distinct \(a,b\in D\) which cannot be joined by a finite path in \(\Gamma(D)\).

\textbf{Unresolved scope qualifications.} The source does not assume \(G=\langle D\rangle\), that \(D\) is a single conjugacy class, or that \(D\) contains every involution of \(G\); none of these conditions is imposed here. The graph depends on the chosen set \(D\), and the source does not explicitly distinguish a classification of groups equipped with a specified \(D\) from a classification of groups admitting some such \(D\). It also does not exclude \(D=\varnothing\) or specify whether the graph with no vertices counts as disconnected. Consequently, that boundary case remains unresolved.
\subsection{Short English statement}
Describe the finite groups having a conjugation-invariant set of involutions whose commuting graph is disconnected. The source leaves the classification's dependence on the chosen set and the empty-set convention unspecified, and does not require the involutions to generate the group.
\subsection{Sources}
\begin{itemize}
\item Local UnsolvedMath registry, supplied record with permanent ID \(700373\), exact source problem code \texttt{catalogue entry}, title \emph{Kourovka Notebook Problem 6.26}; upstream revision \texttt{d371ee917ac9647fbb06cd6fc529deec6f3dda82}.
\item \emph{The Kourovka Notebook (Unsolved Problems in Group Theory)}, original problem 6.26, sixth issue (1978); supplied reference to the twenty-first issue edition (2026), arXiv identifier \texttt{1401.0300v48}. The supplied collection description names E.\ I.\ Khukhro and V.\ D.\ Mazurov as editors.
\item Proposed by A.\ A.\ Makhn\"ev in 1978, according to the supplied attribution. The registry labels the problem \emph{open}; this is a supplied status, not an independently verified mathematical or literature assessment. The cited source has not been independently consulted for this draft.
\end{itemize}
\end{document}
